\section{Predicted Class Correction}
\label{sec:method}

\subsection{Setup}
\label{sec:setup}

Let $(X_i,Y_i)_{i=1}^{n}$ be calibration pairs and $(X_{n+1},Y_{n+1})$ a test pair, all
exchangeable, with labels in $\mathcal{Y}=\{1,\dots,K\}$. A conformal score
$s:\mathcal{X}\times\mathcal{Y}\to\mathbb{R}$ is larger when a pair conforms less; we use
$s(x,y)=1-\hat p(y\mid x)$, the LAC score of \citet{sadinle2019least}, and report APS
\cite{romano2020classification}, RAPS \cite{angelopoulos2021uncertainty} and SAPS
\cite{huang2024conformal} in \cref{tab:depth}. Write $n_y$ for the number of calibration
examples of class $y$, and let
\begin{equation}
\hat q = \mathrm{Quantile}\Big(\{s(X_i,Y_i)\}_{i=1}^n;\ \tfrac{\lceil (n+1)(1-\alpha)\rceil}{n}\Big)
\label{eq:qglobal}
\end{equation}
be the marginal conformal quantile. For a set-valued procedure $\mathcal{C}$ we write
\begin{equation}
\begin{aligned}
\mathrm{CondCov}(\mathcal{C},y) &= \mathbb{P}\bigl(Y\in\mathcal{C}(X)\mid Y=y\bigr), \\
W(\mathcal{C}) &= \min_{y\in\mathcal{Y}} \mathrm{CondCov}(\mathcal{C},y),
\end{aligned}
\label{eq:worst}
\end{equation}
so $W$ is the worst-class coverage this paper reports. Classwise conformal prediction replaces
$\hat q$ by a per-class quantile $\hat q_y$ computed from the $n_y$ examples of class $y$;
equivalently it applies a per-class \emph{offset}
\begin{equation}
\delta_y \;=\; \hat q_y - \hat q ,
\label{eq:delta}
\end{equation}
which is positive for classes that need a looser threshold. When $n_y=0$, $\hat q_y$ and
hence $\delta_y$ are undefined.

\paragraph{Held-out classes.} We study the zero-example case directly. A fraction of classes
is designated as \emph{held out} and all of their calibration rows are deleted, so that
$n_y$ is exactly zero for them. Results are then reported in two tables that
are never averaged together, one for the classes that kept their data and one for the held-out
classes, since a method that wins on the held-out classes by sacrificing the others has not
solved the problem and a single pooled number would conceal that.

\paragraph{Choice of intervention.} A natural alternative is to improve $s(x,y)$ itself, by a
better model, better probability calibration, or aggregation over transformed copies of the
input. Any such change reshapes the score for every class at once, while class-conditional
coverage is set by where each class's threshold falls in its own score distribution. Reducing
the spread of per-class coverage therefore needs a degree of freedom per class, which
$\delta_y$ supplies and a score-level change does not; \cref{app:score-level} measures this.

\paragraph{Matched set size.} Coverage can always be bought with larger sets, so comparing
coverage at unmatched set size compares nothing. Writing
$\mathcal{C}_\delta(x)=\{y: s(x,y)\le \hat q+\delta_y\}$ for the sets induced by an offset
vector $\delta\in\mathbb{R}^{K}$, every method we report is shifted by the scalar
\begin{equation}
\sigma(\delta) \;=\;
\inf\Bigl\{\,\sigma\in\mathbb{R} \;:\;
\tfrac{1}{N}\textstyle\sum_{i=1}^{N}\bigl|\mathcal{C}_{\delta+\sigma}(X_i)\bigr| \;\ge\; T \Bigr\}
\label{eq:sigma}
\end{equation}
for a common target size $T$, and worst-class coverage is read at $\delta+\sigma(\delta)$. We
call $\bar\delta = \delta + \sigma(\delta)$ the \emph{matched} offset, and every
comparison in \cref{sec:exp} is read at it.

\subsection{Calibration-label-free class descriptors}
\label{sec:phi}

For a linear classifier head with weight rows $w_y\in\mathbb{R}^d$ and biases $b_y$, we form a
descriptor vector $\varphi(y)\in\mathbb{R}^{p}$, $p=9$, from
\begin{enumerate}[itemsep=1pt,topsep=2pt,leftmargin=*]
\item $\|w_y\|_2$ and $b_y$, the scale and offset of the class logit;
\item cosine distances from $w_y$ to its $k$ nearest competing rows for
      $k\in\{5,10,50\}$, the distance to the nearest row, the spread of that
      neighbourhood, and the cosine distance to the mean row, all of which measure how
      crowded the class is;
\item $\log$ of the class prevalence in the training set.
\end{enumerate}
The eight geometric features are functions of the model's parameters alone, so they are
defined for a class with no examples of any kind and carry no sampling noise. Prevalence is
counted from the training labels, which are available whenever the model is; the descriptors
are free of \emph{calibration} labels, though not of labels altogether. No result depends
on it: removing prevalence costs $0.011$ in worst-class coverage without reaching significance
($p=0.30$), while prevalence alone is significantly worse than the full set
(\cref{tab:ablation}). The
feature that most closely resembles the regression target is held out, so no reported number
can be attributed to it.

\subsection{The correction}
\label{sec:correction}

\paragraph{Step 1: regress.} On the classes with at least $n_{\mathrm{cal}}$ calibration
examples, a count fixed in advance, we observe $\delta_y$ from
\cref{eq:delta} and fit a ridge regression $g_\theta:\mathbb{R}^p\to\mathbb{R}$ to predict it
from $\varphi(y)$. There is no network and no training loop: $g_\theta$ is a closed-form fit
over at most $K$ points with $p\le 15$ features, which costs milliseconds
(\cref{tab:runtime}).

\paragraph{Step 2: shrink.} The raw prediction is applied with a scalar $\lambda$,
\begin{equation}
\tilde\delta_y \;=\; \lambda\cdot g_\theta(\varphi(y)) ,
\label{eq:shrink}
\end{equation}
with $\lambda$ chosen over a fixed grid $\Lambda$ by maximising worst-class coverage at
matched size on the labelled classes only,
\begin{equation}
\hat\lambda \;=\; \operatorname*{arg\,max}_{\lambda\in\Lambda}\;
\widehat W_{\mathcal{Y}_{\mathrm{lab}}}\bigl(\mathcal{C}_{\overline{\lambda g_\theta(\varphi)}}\bigr),
\label{eq:lambda}
\end{equation}
where $\widehat W_{\mathcal{Y}_{\mathrm{lab}}}$ is the empirical version of \cref{eq:worst}
restricted to $\mathcal{Y}_{\mathrm{lab}}$. The restriction matters, because choosing
$\lambda$ with help from the held-out classes would let a free parameter manufacture the very
result the held-out table is meant to test, and it would void \cref{prop:marginal}.

Shrinkage is what makes the method usable, and \cref{prop:sup} says why;
\cref{fig:method} shows the fit it is applied to. Worst-class
coverage is controlled by the \emph{largest} error in the offsets, a quantity mean squared
error does not bound, so a predictor with a respectable $R^2$ can still be harmful at full
strength: $\lambda=1$ costs $0.584$ relative to the selected $\lambda\approx0.1$, while
$\lambda=0$ recovers the marginal threshold exactly (\cref{tab:ablation}).

\begin{proposition}[Worst-class coverage is a sup-norm functional]
\label{prop:sup}
Let $F_y(t)=\mathbb{P}(s(X,Y)\le t\mid Y=y)$ and suppose each $F_y$ is $L$-Lipschitz. For any
two offset vectors $\delta,\delta'\in\mathbb{R}^{K}$ with matched versions
$\bar\delta,\bar\delta'$ from \cref{eq:sigma},
\begin{equation}
\bigl|\,W(\mathcal{C}_{\bar\delta}) - W(\mathcal{C}_{\bar\delta'})\,\bigr|
\;\le\; L\,\bigl\|\bar\delta-\bar\delta'\bigr\|_{\infty}.
\label{eq:supbound}
\end{equation}
\end{proposition}

\begin{proof}
$\mathrm{CondCov}(\mathcal{C}_{\bar\delta},y)=F_y(\hat q+\bar\delta_y)$, so by Lipschitzness every $y$ obeys
\begin{equation}
|F_y(\hat q+\bar\delta_y)-F_y(\hat q+\bar\delta'_y)|
\;\le\; L\,\|\bar\delta-\bar\delta'\|_\infty .
\end{equation} Two functions of $y$ that agree
pointwise to within $\varepsilon$ have minima that agree to within $\varepsilon$, and
\cref{eq:worst} is such a minimum.
\end{proof}

Taking $\delta'$ to be the oracle offsets bounds our worst-class coverage below by the
oracle's, less $L$ times the sup-norm error. The bound is in the sup norm, and a ridge fit
controls a mean, so a single badly predicted class can cost as much as many well predicted ones
save. Shrinking interpolates between two sup-norm errors: at $\lambda=0$ the error is
$\|\bar\delta^\star\|_\infty$, the whole correction foregone, and at $\lambda=1$ it is the
largest prediction error of $g_\theta$. When the second exceeds the first, an interior
$\lambda$ beats both endpoints, which is what \cref{tab:ablation} measures.

\paragraph{Accuracy required of $g_\theta$.} We have no theory saying when $\varphi$ predicts
$\delta$ well, and the method works exactly as far as the regression does; a quadratic or
kernel head does not change that (\cref{app:heads}). Across our three
datasets the two line up monotonically. The cross-validated error of $g_\theta$ is $0.036$ on
ImageNet, $0.072$ on Pl@ntNet-300K and $0.127$ on iNaturalist-2018, against worst-class
improvements of $0.058$, $0.019$ and $0.000$. Because \cref{eq:lambda} can return $\hat\lambda=0$, a poor regression is
abstained from, which is what happens on iNaturalist-2018 in every split (\cref{sec:depth}).

\paragraph{Step 3: recalibrate.} Finally one scalar $c$ is refit, on labelled-class rows only,
so that marginal coverage returns to $1-\alpha$. The prediction set is
\begin{equation}
\mathcal{C}(x) \;=\; \bigl\{\, y : s(x,y) \le \hat q + \tilde\delta_y + c \,\bigr\}.
\label{eq:set}
\end{equation}

\begin{algorithm}[t]
\small
\caption{Predicted Class Correction}
\label{alg:pcc}
\begin{algorithmic}[1]
\Require score $s$, level $\alpha$, target size $T$, grid $\Lambda$, count $n_{\mathrm{cal}}$
\Require disjoint calibration rows $\mathcal{I}_{\mathrm{fit}},\mathcal{I}_{\mathrm{rec}}$;
         labelled classes $\mathcal{Y}_{\mathrm{lab}}$; descriptors $\varphi$
\State $\hat q \gets$ marginal conformal quantile \eqref{eq:qglobal} on
       $\mathcal{I}_{\mathrm{fit}}$
\For{$y\in\mathcal{Y}_{\mathrm{lab}}$ with $n_y\ge n_{\mathrm{cal}}$}
  \State $\delta_y \gets \hat q_y - \hat q$ \Comment{observed offset \eqref{eq:delta}}
\EndFor
\State $g_\theta \gets$ ridge fit of $\delta$ on $\varphi$ over those classes
\State $\hat\lambda \gets$ \eqref{eq:lambda}, on $\mathcal{Y}_{\mathrm{lab}}$ only
\State $\tilde\delta_y \gets \hat\lambda\, g_\theta(\varphi(y))$ for \emph{every}
       $y\in\mathcal{Y}$ \Comment{including $n_y=0$}
\State $c \gets$ conformal quantile of
       $\{s(X_i,Y_i)-\tilde\delta_{Y_i}\}_{i\in\mathcal{I}_{\mathrm{rec}}}$
\State \Return $\mathcal{C}(x)=\{y : s(x,y)\le \hat q+\tilde\delta_y+c\}$
\end{algorithmic}
\end{algorithm}

\subsection{Data splitting}
\label{sec:splits}

Every quantity above is estimated, so it matters which rows and which classes each one may
see. Classes are partitioned into labelled and held-out, and independently the rows of every
class into a calibration and an evaluation part; \cref{tab:splits} in the appendix states the
combination each quantity is fitted on. Two invariants hold throughout: nothing is fitted on evaluation rows,
and nothing is fitted on a held-out class. The numbers reported for held-out classes are
therefore free of selection effects from $\lambda$, $n_{\mathrm{cal}}$ and $c$, while those for
labelled classes are not, which is the second reason the two tables are kept apart.


\subsection{Coverage guarantee}
\label{sec:guarantee}

Let $\mathcal{I}_{\mathrm{fit}}$ and $\mathcal{I}_{\mathrm{rec}}$ be disjoint sets of
calibration rows, of sizes $n_{\mathrm{fit}}$ and $n_{\mathrm{rec}}$.

\begin{proposition}[Marginal coverage is preserved]
\label{prop:marginal}
Let $(X_i,Y_i)_{i\in\mathcal{I}_{\mathrm{rec}}}$ and $(X_{n+1},Y_{n+1})$ be exchangeable, and
let $\tilde\delta_\cdot$ be any function of $\varphi(\cdot)$ and of
$(X_i,Y_i)_{i\in\mathcal{I}_{\mathrm{fit}}}$. Let $c$ be the conformal quantile at level
$\lceil (n_{\mathrm{rec}}+1)(1-\alpha)\rceil/n_{\mathrm{rec}}$ of the shifted scores
$\{s(X_i,Y_i)-\tilde\delta_{Y_i}\}_{i\in\mathcal{I}_{\mathrm{rec}}}$. Then
$\mathbb{P}\bigl(Y_{n+1}\in\mathcal{C}(X_{n+1})\bigr)\ge 1-\alpha$.
\end{proposition}

\begin{proof}[Proof sketch]
Define $\tilde s(x,y)=s(x,y)-\tilde\delta_y$. Since $\tilde\delta_\cdot$ is a function of
$\varphi(\cdot)$ alone, and $\varphi$ is computed from the model's weights and does not involve
these calibration rows, $\tilde s$ is a \emph{fixed} score function with respect to the
sample. Exchangeability is closed under applying a fixed map pointwise, so
$\tilde s(X_1,Y_1),\dots,\tilde s(X_{n+1},Y_{n+1})$ are exchangeable and the standard split
conformal argument applies verbatim to $\tilde s$. Membership in $\mathcal{C}$ is exactly the
event $\tilde s(X_{n+1},Y_{n+1})\le c$ up to the constant shift. Full proof in
\cref{app:proof}.
\end{proof}

The proposition holds because only \emph{one} degree of freedom is calibrated on
$\mathcal{I}_{\mathrm{rec}}$; $\tilde\delta_y$ may be arbitrarily complex, since it merely
redefines the score.

\paragraph{Reuse of calibration rows.} The premise names two disjoint row sets, and the
configuration used for every table in \cref{sec:exp} does not provide them: $g_\theta$,
$\lambda$ and $c$ are all estimated on the same calibration rows, so $\tilde\delta$ and $c$ are
dependent and the proposition does not apply as stated. On ten ImageNet splits, at the thresholds as they are
emitted, \textsc{Pcc} covers $0.8994$ where the marginal threshold covers $0.9012$ on the same
rows, so the reuse costs $0.0018$ of marginal coverage, roughly twenty times the finite-sample
slack $1/(n+1)$ at our calibration sizes. \textsc{Pcc}-split holds a quarter of the calibration
rows back for $c$ alone, which satisfies the premise, and it removes that deficit entirely
while changing worst-class coverage by $+0.005$ with an interval containing zero
(\cref{app:split}). The tables below use the reuse configuration because it is what every
earlier run used; the split variant is the one we would recommend deploying.

\paragraph{Scope of the claim.} No method can carry a finite-sample class-conditional
guarantee at $n_y=0$, since there is nothing to condition on. Our claim for those classes is
empirical, namely measured worst-class coverage at matched set size. And $\lambda$ and
$n_{\mathrm{cal}}$ are chosen on the labelled label space, which makes the labelled-class
table optimistic; this is why the two tables are kept apart.
